\section{Conclusion}
\label{sec:conclusion}

We adapted ViT-B/16 to a 10,000-class fine-grained benchmark with 45 images per class within a budget of 96\,GB$\cdot$h per run.
The largest gain came from the least glamorous component.
Initializing the 10,000-way head with imprinted class prototypes instead of a linear probe fitted on the training features improved accuracy by 2.5--3.1 points, for both full fine-tuning and LoRA.
With this head, full fine-tuning reaches 76.22\% on the CodaBench validation set in 1.03 GPU-hours.
LoRA on all linear layers reaches 71.90\% while training 1.375\% of the backbone.
After merging, LoRA has zero inference overhead.
For PEFT under this budget, three choices mattered:
\begin{itemize}[leftmargin=*,itemsep=0pt,topsep=2pt]
    \item where the adaptation capacity goes (all linear layers, not only attention);
    \item whether the update can be merged (LoRA \vs AdaptFormer);
    \item a short schedule (about 3 epochs).
\end{itemize}

\paragraph{Limitations.}
Every configuration was trained once (seed 0), and the holdout split has a noise level of about $\pm$0.3 points, so small differences such as the 0.26-point gain from the two extra images per class should be read with care.
The dataset is anonymized, so we cannot relate our findings to its domain.
One ablation batch ran on a different GPU; its accuracies are comparable, as checked by repeating one run, but its speeds and latencies are not.
Our explanation of why the linear-probe head hurts is consistent with the loss curves but is not a controlled proof.
We did not explore higher input resolution, other PEFT families such as VPT or BitFit, or larger budgets.
\todo{Phase 2: report the CodaBench test accuracy of the final models.}
\todo{Check the course page limit (the CVPR limit is 8 pages excluding references).}
